\chapter{Cơ sở lý thuyết}\label{chap:theory}
Chương này trình bày các cơ sở lý thuyết mà đồ án sử dụng. Phần đầu nêu khái niệm hệ thống gia sư thông minh, phương pháp Socratic, giàn giáo, tự giải thích và lỗi quan niệm, tức là những cơ sở sư phạm quyết định hành vi của gia sư. Phần tiếp theo giới thiệu mô hình ngôn ngữ lớn, các mô hình và hạ tầng suy luận được dùng cùng những giới hạn của chúng. Các phần sau trình bày truy xuất tăng cường, đồ thị tri thức, mô hình người học và truy vết tri thức, bandit ngữ cảnh cùng đánh giá ngoài chính sách, xử lý tiếng nói, phương pháp đánh giá và các kỹ thuật phần mềm nền. Mỗi phần kết thúc bằng điểm mà đồ án kế thừa hoặc điều chỉnh, để Chương~\ref{chap:design} có thể dẫn lại trực tiếp.

\section{Hệ thống gia sư thông minh và phương pháp Socratic}

\subsection{Hệ thống gia sư thông minh}
Hệ thống gia sư thông minh (ITS) là phần mềm dạy học mô phỏng một gia sư người: hệ thống theo dõi từng bước làm của học sinh, phản hồi và gợi ý đúng lúc, rồi điều chỉnh theo người học. Kiến trúc kinh điển của ITS gồm bốn mô-đun. Mô-đun miền tri thức chứa kiến thức cần dạy, lời giải và các lỗi thường gặp. Mô-đun người học ước lượng học sinh đã biết gì và đang sai gì. Mô-đun sư phạm quyết định dạy gì tiếp theo và gợi ý ra sao. Mô-đun giao diện là nơi học sinh tương tác với hệ thống. Hình~\ref{fig:its-modules} minh họa bốn mô-đun này cùng vị trí mà đồ án đặt mô hình ngôn ngữ.

\begin{figure}[ht]
\centering
\begin{tikzpicture}[font=\footnotesize,
  mod/.style={draw=black,line width=0.5pt,fill=white,align=center,minimum height=17mm,text width=42mm,inner sep=2pt},
  llm/.style={mod,fill=gray!30,dashed,line width=0.7pt,text width=40mm,minimum height=13mm},
  arr/.style={-{Stealth[length=2mm,width=1.6mm]},line width=0.6pt},
  lab/.style={font=\scriptsize\itshape,fill=white,inner sep=1pt}]
\node[mod] (dom) at (0,0) {\textbf{Mô-đun miền tri thức}\\\scriptsize kiến thức, lời giải, lỗi thường gặp\\[1pt]\scriptsize\itshape CAVAT: gói chương, đồ thị};
\node[mod] (stu) at (6.0,0) {\textbf{Mô-đun người học}\\\scriptsize ước lượng đã biết gì, sai gì\\[1pt]\scriptsize\itshape CAVAT: phép chiếu mức hiểu bài};
\node[mod] (tut) at (3.0,-3.3) {\textbf{Mô-đun sư phạm}\\\scriptsize dạy gì tiếp, gợi ý ra sao\\[1pt]\scriptsize\itshape CAVAT: luật N-*, R-*};
\node[mod] (ui) at (3.0,-6.6) {\textbf{Mô-đun giao diện}\\\scriptsize nơi học sinh tương tác\\[1pt]\scriptsize\itshape CAVAT: web, giọng nói};
\node[llm] (lm) at (9.4,-4.95) {\textbf{Mô hình ngôn ngữ}\\\scriptsize chỉ diễn đạt lời gia sư};
\draw[arr] (dom.south) -- ++(0,-0.6) -| ([xshift=-6mm]tut.north);
\draw[arr] (stu.south) -- ++(0,-0.6) -| ([xshift=6mm]tut.north);
\draw[arr] (tut) -- (ui);
\draw[arr,dashed] (tut.east) -| (lm.north);
\draw[arr,dashed] (lm.south) |- ([yshift=3mm]ui.east);
\draw[arr] ([yshift=-3mm]ui.east) -- (12.4,-6.9) |- node[lab,pos=0.2,right]{câu trả lời} (stu.east);
\end{tikzpicture}

\caption{Bốn mô-đun của một hệ thống gia sư thông minh và cách CAVAT hiện thực từng mô-đun. Mô hình ngôn ngữ nằm giữa mô-đun sư phạm và giao diện, chỉ với vai trò diễn đạt.}
\label{fig:its-modules}
\end{figure}

VanLehn phân biệt các hệ thống gia sư theo độ mịn của tương tác: hỗ trợ theo câu trả lời cuối, theo từng bước, hoặc theo từng ý trong một bước~\cite{vanlehn2011}. Hỗ trợ càng mịn thì hiệu quả càng gần với gia sư người. Cognitive Tutor và AutoTutor là hai hệ tiêu biểu. AutoTutor dẫn dắt học sinh bằng hội thoại có cấu trúc gồm các nước đi gợi ý, thúc đẩy và sửa sai~\cite{graesser2005}. Đồ án giữ nguyên bốn mô-đun cổ điển và hỗ trợ ở mức từng lượt nói của học sinh. Điểm khác biệt nằm ở chỗ mô hình ngôn ngữ được đặt vào giữa mô-đun sư phạm và giao diện, và quyết định sư phạm không rời khỏi mô-đun sư phạm.

\subsection{Phương pháp Socratic}\label{sec:socratic-ref}
Phương pháp Socratic dạy bằng câu hỏi dẫn người học tự phát hiện điều cần biết thay vì kể cho họ nghe. Nguyên tắc này khó kiểm chứng nếu chỉ phát biểu ở mức triết lý, vì một lời giảng có thể vừa hỏi vừa để lộ đáp án. Đồ án vì vậy tách phương pháp Socratic thành sáu ràng buộc hành vi mà bộ kiểm chứng có thể kiểm tra được. Gia sư không nói ra đáp án cuối trước khi học sinh tự giải được, trừ một nước đi được phép là giải thích đầy đủ sau khi cùng làm từng bước đã thất bại. Mỗi lượt chỉ chứa một câu hỏi và câu hỏi đó đặt ở cuối lượt. Gia sư hỏi cách làm và lý do trước khi sửa. Mức trợ giúp bắt đầu từ bậc nhẹ nhất rồi tăng dần, và gia sư rút trợ giúp lại khi học sinh đã làm được một bước. Những lượt không phải một lần thử, chẳng hạn lạc đề, hỏi lại hay phàn nàn, không bị tính là sai.

Các ràng buộc này cũng là các tiêu chí đo ở Chương~\ref{chap:eval}. Hai tiêu chí đầu, không lộ đáp án và một câu hỏi mỗi lượt, đo được hoàn toàn bằng luật. Bốn tiêu chí còn lại phụ thuộc vào trạng thái của buổi học nên được thực thi ở chính sách dạy thay vì ở bộ kiểm chứng.

\section{Giàn giáo, tự giải thích và lỗi quan niệm}

\subsection{Giàn giáo và nghịch lý hỗ trợ}
Giàn giáo là sự hỗ trợ tạm thời giúp người học làm được điều mình chưa tự làm được, và sự hỗ trợ đó được rút dần khi người học làm chủ nhiệm vụ~\cite{wood1976}. Khái niệm này gắn với vùng phát triển gần của Vygotsky, tức khoảng giữa điều học sinh tự làm được và điều học sinh làm được khi có người giúp~\cite{vygotsky1978}. Hình~\ref{fig:zpd} minh họa mối quan hệ đó. Gia sư hiệu quả luôn nhắm vào vùng giữa: nhiệm vụ nằm trong vùng tự làm được thì không cần giúp, còn nhiệm vụ nằm ngoài vùng phát triển gần thì sự giúp đỡ không đủ để học sinh tiến bộ.

\begin{figure}[ht]
\centering
\begin{tikzpicture}[font=\footnotesize]
\fill[gray!8] (0,0) ellipse (5.0cm and 3.1cm);
\fill[gray!22] (0,0) ellipse (3.4cm and 2.2cm);
\fill[gray!42] (0,0) ellipse (1.6cm and 0.8cm);
\draw[line width=0.6pt] (0,0) ellipse (5.0cm and 3.1cm);
\draw[line width=0.6pt] (0,0) ellipse (3.4cm and 2.2cm);
\draw[line width=0.6pt] (0,0) ellipse (1.6cm and 0.8cm);
\node[align=center] at (0,0) {\textbf{Tự làm được}};
\node[align=center] at (0,1.5) {\textbf{Làm được khi có hỗ trợ}\\\scriptsize\itshape vùng phát triển gần};
\node[align=center] at (0,-2.65) {\textbf{Chưa làm được}};
\draw[-{Stealth[length=2mm,width=1.6mm]},line width=0.7pt] (-7.3,0.9) -- (-1.9,0.3);
\node[font=\scriptsize\itshape,anchor=south west] at (-7.5,1.0) {giàn giáo hỗ trợ};
\draw[-{Stealth[length=2mm,width=1.6mm]},line width=0.7pt] (7.3,-0.9) -- (3.5,-0.9);
\node[font=\scriptsize\itshape,anchor=north east] at (7.5,-1.0) {rút dần hỗ trợ};
\end{tikzpicture}

\caption{Vùng phát triển gần. Giàn giáo đưa nhiệm vụ từ vùng làm được khi có hỗ trợ vào vùng tự làm được, và hỗ trợ được rút dần khi năng lực tăng.}
\label{fig:zpd}
\end{figure}

Koedinger và Aleven gọi mâu thuẫn trong việc quyết định mức hỗ trợ là nghịch lý hỗ trợ: giúp quá nhiều làm học sinh mất cơ hội tự nghĩ, giúp quá ít làm học sinh bế tắc~\cite{koedinger2007}. Nghiên cứu về hành vi xin trợ giúp cho thấy học sinh thường lạm dụng gợi ý hoặc né tránh gợi ý, nên hệ thống cần quy tắc rõ ràng về thời điểm và mức độ hỗ trợ~\cite{aleven2016}. Các hệ gia sư thường dùng thang gợi ý đi từ chung đến cụ thể, với bậc cuối là chính lời giải. Đồ án dùng một thang tám bậc từ H0 đến H7 theo hai trục là bậc (giúp bao nhiêu) và mục tiêu (giúp vào đâu). Mỗi bậc có một trọng số độc lập do khung đánh giá quy định. Gia sư và bộ chấm dùng chung các trọng số đó, nên hai bên cùng một thước đo về mức học sinh đã được giúp. Chương~\ref{chap:design} mô tả chi tiết thang này.

\subsection{Tự giải thích và hiệu ứng đảo ngược chuyên môn}
Hiệu ứng tự giải thích cho thấy học sinh tự giải thích các bước và lý do của mình thì hiểu sâu và nhớ lâu hơn~\cite{chi1989,chi1994}. Chi và cộng sự phân tích việc học từ gia sư như quá trình học sinh tự xây dựng kiến thức chứ không phải tiếp nhận thụ động~\cite{chi2001}. Kết quả này là cơ sở của hai nước đi trong đồ án, một nước đi hỏi ``vì sao'' và một nước đi mời học sinh giảng lại cách làm.

Hiệu ứng đảo ngược chuyên môn đặt giới hạn cho việc áp dụng ở trên. Một kiểu hỗ trợ có ích với người mới có thể không giúp được, thậm chí gây cản trở, với người đã thành thạo~\cite{kalyuga2003}. Gia sư vì vậy chỉ hỏi ``vì sao'' khi kiến thức còn mới hoặc chưa vững, và hỏi ít đi khi học sinh đã vững.

\subsection{Lỗi quan niệm và xung đột nhận thức}
Lỗi quan niệm là một cách hiểu sai có hệ thống chứ không phải sai ngẫu nhiên. Resnick và cộng sự ghi nhận ở số thập phân các lỗi điển hình như cho rằng số nào dài hơn thì lớn hơn (viết $0{,}25>0{,}3$), do học sinh chuyển cách nghĩ của số tự nhiên sang~\cite{resnick1989}. Một lỗi quen thuộc khác là nhầm hàng phần mười với hàng phần trăm, chẳng hạn viết $9/10$ thành $0{,}09$.

Xung đột nhận thức bắt nguồn từ lý thuyết thay đổi khái niệm: người dạy đưa cho học sinh một phản ví dụ để học sinh tự thấy cách nghĩ của mình dẫn tới mâu thuẫn, thay vì nói thẳng rằng cách nghĩ đó sai~\cite{posner1982}. Đồ án gắn mỗi lỗi quan niệm mà giáo viên khai báo với một phản ví dụ và dùng phản ví dụ đó như một nước đi sư phạm riêng. Bộ nhận diện lỗi quan niệm là một luật tất định dựa trên dạng sai đặc trưng, không dựa trên phán đoán của mô hình ngôn ngữ.

\subsection{Khen nỗ lực}
Mueller và Dweck cho thấy khen trí thông minh làm trẻ ngại thử thách và nản khi gặp khó, còn khen nỗ lực thì cho kết quả ngược lại~\cite{mueller1998}. Bộ kiểm chứng của đồ án vì thế chặn lời khen trí thông minh. Bộ kiểm chứng cũng chặn lời khen đứng trước một câu trả lời sai, vì lời khen sai chỗ làm học sinh hiểu nhầm rằng câu trả lời đã đúng.

\section{Mô hình ngôn ngữ lớn}

\subsection{Nguyên lý hoạt động}
Mô hình ngôn ngữ ước lượng xác suất của một chuỗi từ. Các mô hình ngôn ngữ lớn hiện nay là mô hình tự hồi quy: mô hình nhận một đoạn văn bản, dự đoán token tiếp theo rồi lặp lại bước đó để sinh văn bản. Văn bản được cắt thành các đơn vị con bằng thuật toán như Byte-Pair Encoding, và tiếng Việt có dấu thường tốn nhiều token hơn tiếng Anh cho cùng một ý. Chi phí và giới hạn độ dài vì vậy bị ảnh hưởng trực tiếp khi dùng tiếng Việt.

Kiến trúc Transformer là nền của các mô hình này~\cite{vaswani2017}. Cơ chế tự chú ý cho phép mỗi token ``nhìn'' mọi token đứng trước với trọng số học được, nhiều đầu chú ý chạy song song và mã hóa vị trí cho biết thứ tự. Các mô hình sinh văn bản dùng kiến trúc chỉ có bộ giải mã. Giai đoạn tiền huấn luyện dạy mô hình dự đoán token tiếp theo trên khối văn bản rất lớn, nhờ đó mô hình học được ngôn ngữ và kiến thức chung. Giai đoạn tinh chỉnh theo chỉ dẫn và học tăng cường từ phản hồi của người sau đó dạy mô hình làm theo yêu cầu và trả lời theo cách người đánh giá ưa thích~\cite{ouyang2022}.

\subsection{Cách sử dụng trong hệ thống}
Mô hình làm theo chỉ dẫn và ví dụ đặt trong lời nhắc mà không cần huấn luyện lại, một khả năng gọi là học trong ngữ cảnh~\cite{brown2020}. Lời nhắc hệ thống đặt vai trò và luật cho mô hình, còn cửa sổ ngữ cảnh giới hạn tổng độ dài của lời nhắc và câu trả lời. Một số mô hình có chế độ suy luận, trong đó mô hình sinh thêm một chuỗi suy nghĩ trước câu trả lời; chế độ này chính xác hơn với bài khó nhưng tốn nhiều token và chậm hơn. Chương trình có thể yêu cầu mô hình trả về JSON theo một lược đồ để đọc được kết quả bằng máy. Tham số nhiệt độ điều khiển độ ngẫu nhiên của đầu ra, và nhiệt độ thấp cho câu trả lời ổn định hơn.

\subsection{Các giới hạn liên quan tới dạy học}
Năm giới hạn của mô hình ngôn ngữ liên quan trực tiếp tới dạy học. Ảo giác là hiện tượng mô hình nói tự tin những điều không đúng hoặc không có căn cứ~\cite{ji2023}. Xu hướng chiều người dùng làm mô hình dễ đồng ý với học sinh, kể cả khi học sinh sai. Việc được tối ưu để hữu ích khiến mô hình hay giải hộ và để lộ đáp án~\cite{macina2023}. Tấn công chèn lời nhắc xảy ra khi lời của người dùng hoặc dữ liệu ngoài chứa chỉ thị làm mô hình đổi hành vi~\cite{perez2022,greshake2023}. Tính không tất định và khó kiểm toán là giới hạn cuối: cùng một đầu vào có thể cho đầu ra khác nhau, và mô hình không giải thích được vì sao chọn một cách dạy. Bảng~\ref{tab:llm-limits} đối chiếu từng giới hạn với cách đồ án ứng phó.

\begin{table}[ht]
\centering
\caption{Giới hạn của mô hình ngôn ngữ lớn trong dạy học và cơ chế ứng phó của đồ án.}
\label{tab:llm-limits}
\small
\begin{tabularx}{\linewidth}{@{}p{3.2cm}XX@{}}
\toprule
Giới hạn & Biểu hiện trong dạy học & Cơ chế ứng phó \\
\midrule
Ảo giác & Khẳng định kiến thức hoặc số liệu không có căn cứ & Truy xuất có ràng buộc theo đồ thị; kiểm số và thuật ngữ so với nguồn \\
Chiều người dùng & Đồng ý với đáp án sai của học sinh & Bộ chấm tất định quyết định đúng sai \\
Lộ đáp án & Giải hộ khi học sinh nhờ & Đáp án câu kiểm tra không có trong lời nhắc; bộ kiểm chứng quét đáp án \\
Chèn lời nhắc & Lời học sinh làm đổi vai và luật & Lời học sinh và đoạn nguồn chỉ là dữ liệu; chính sách nằm ngoài mô hình \\
Không tất định & Mức hiểu bài khó tái lập & Mô hình không ghi điểm; mức hiểu bài là phép chiếu từ sổ sự kiện \\
\bottomrule
\end{tabularx}
\end{table}

Các giới hạn này là lý do đồ án không để mô hình ngôn ngữ quyết định nước đi, chấm điểm hay ghi mức hiểu bài. Phần còn lại của chương xem xét các thành phần tất định thay thế cho những việc đó.

\section{Các mô hình và hạ tầng suy luận được sử dụng}

\subsection{Họ mô hình Qwen}
Qwen là họ mô hình ngôn ngữ lớn do đội Qwen của Alibaba Cloud phát triển, với các thế hệ chính gồm Qwen~\cite{qwen2023}, Qwen2.5~\cite{qwen25} và Qwen3~\cite{qwen3}. Phần lớn mô hình được công bố trọng số mở, nhiều phiên bản theo giấy phép Apache 2.0. Các thế hệ gần đây dùng Transformer chỉ giải mã với chú ý truy vấn nhóm để giảm bộ nhớ khi suy luận, mã hóa vị trí xoay và chuẩn hóa RMSNorm. Qwen3 được tiền huấn luyện trên khoảng 36 nghìn tỉ token, hỗ trợ khoảng 119 ngôn ngữ và phương ngữ trong đó có tiếng Việt, gồm cả mô hình dày và mô hình hỗn hợp chuyên gia từ 0,6 đến 235 tỉ tham số. Một mô hình Qwen3 gộp hai chế độ, chế độ suy nghĩ cho bài khó và chế độ không suy nghĩ cho phản hồi nhanh~\cite{qwen3}.

Đồ án dùng phiên bản \texttt{qwen/qwen3.8-27b} phục vụ qua Groq làm bộ não của gia sư ở giai đoạn đầu. Thẻ mô hình của Groq mô tả đây là mô hình dày 27 tỉ tham số, đa phương thức, có 64 lớp với kiến trúc lai Gated DeltaNet và Gated Attention, cửa sổ ngữ cảnh 131\,072 token và hai chế độ suy nghĩ chuyển bằng tham số \texttt{reasoning\_effort}~\cite{groqdocs}. Bốn lý do dẫn tới lựa chọn này. Tiếng Việt của mô hình tốt trong các mô hình miễn phí đã thử. Mô hình miễn phí trên Groq với hạn mức đủ cho giai đoạn thử nghiệm. Mô hình hỗ trợ đầu ra JSON theo lược đồ chặt, điều kiện cần cho bộ chọn nước đi và bộ đọc ngữ cảnh. Khả năng tắt chế độ suy luận là lý do cuối: đo thực tế cho thấy một lượt trả lời ngắn tốn khoảng 392 token khi để mặc định và khoảng 148 token khi tắt suy luận, vì phần suy luận vẫn bị tính vào hạn mức.

\subsection{Các mô hình dự phòng và vai trò của Gemini}
Hai mô hình gpt-oss-120b và gpt-oss-20b của OpenAI được dùng làm dự phòng khi Qwen không dùng được. Cả hai có kiến trúc hỗn hợp chuyên gia và giấy phép Apache 2.0, với 117 và 21 tỉ tham số, trong đó 5,1 và 3,6 tỉ tham số hoạt động cho mỗi token~\cite{gptoss2025}. Mô hình gpt-oss-120b còn là giám khảo trong các thực nghiệm, vì giám khảo cần khác họ với mô hình gia sư được chấm.

Gemini là họ mô hình đa phương thức đọc được văn bản, ảnh và âm thanh~\cite{gemini2023}. Đồ án dùng Gemini cho hai việc: đọc ảnh và nhận dạng ký tự trên các trang tài liệu chụp hoặc quét cùng hình học trong đề, và tạo vector nhúng bằng mô hình \texttt{gemini-embedding-001} với 768 chiều. Chủ dự án cho phép đặt cả Gemini và các mô hình khác vào chuỗi viết lời của gia sư ở giai đoạn sau. Chuỗi hiện hành, theo thứ tự ưu tiên do chủ dự án chốt, là Gemini, gpt-oss-120b, Qwen và gpt-oss-20b, và Chương~\ref{chap:design} trình bày cách cổng AI điều phối chuỗi đó.

Vector nhúng là cách biểu diễn một đoạn văn bản bằng một vector số sao cho các đoạn gần nghĩa có vector gần nhau, và độ gần thường đo bằng cosin của góc giữa hai vector. Mô hình nhúng được huấn luyện bằng học tương phản trên các cặp câu liên quan.

\subsection{Hạ tầng suy luận Groq và hạn mức miễn phí}
Groq cung cấp dịch vụ suy luận trên chip LPU do hãng tự thiết kế, cho tốc độ sinh token cao và độ trễ thấp, với giao diện tương thích OpenAI. Đồ án dùng gói miễn phí của Groq, và gói này có hạn mức theo phút và theo ngày như Bảng~\ref{tab:quota}.

\begin{table}[ht]
\centering
\caption{Hạn mức gói miễn phí của Groq tại thời điểm đo (docs/poc/free-tier-quota.md).}
\label{tab:quota}
\small
\begin{tabularx}{\linewidth}{@{}lXXXX@{}}
\toprule
Mô hình & Lần gọi/phút & Lần gọi/ngày & Token/phút & Token/ngày \\
\midrule
qwen/qwen3.8-27b & 30 & 1\,000 & 8\,000 & 200\,000 \\
gpt-oss-120b, gpt-oss-20b & 30 & 1\,000 & 8\,000 & 200\,000 \\
whisper-large-v3 & 20 & 2\,000 & \multicolumn{2}{l}{7\,200 giây âm thanh/giờ} \\
\bottomrule
\end{tabularx}
\end{table}

Hạn mức theo phút đo được thấp hơn con số công bố: Groq từ chối khi số token đầu vào vượt khoảng 7\,000 trong một phút. Một ngày đủ cho khoảng bốn buổi học dùng hết trần gọi, vì mỗi lần gọi Qwen tốn trung bình 1\,126 token (đo trên hai buổi mô phỏng). Hệ quả thiết kế là mỗi lời nhắc không quá 3\,000 token, mỗi buổi học có trần số lần gọi mô hình, và hệ thống phải xử lý được tình huống mô hình tạm hết hạn mức mà không làm gián đoạn buổi học. Đây là một ràng buộc thiết kế chứ không phải chi tiết triển khai.

\section{Truy xuất thông tin và sinh văn bản tăng cường truy xuất}

\subsection{Khái niệm RAG}
Sinh văn bản tăng cường truy xuất (RAG) là kỹ thuật tìm các đoạn tài liệu liên quan rồi đưa vào lời nhắc trước khi mô hình trả lời, để câu trả lời dựa trên nguồn thay vì dựa vào trí nhớ của mô hình~\cite{lewis2020}. Gao và cộng sự chia một hệ RAG thành các khâu lập chỉ mục, truy xuất, xếp hạng lại và sinh có căn cứ~\cite{gao2023survey}. RAG có hai tác dụng đối với gia sư. Tác dụng thứ nhất là neo lời giảng vào tài liệu của giáo viên. Tác dụng thứ hai là giới hạn phạm vi, vì gia sư không được trả lời bằng nội dung mà học sinh chưa học.

\subsection{Truy xuất thưa theo từ khóa và truy xuất dày theo ngữ nghĩa}
BM25 là hàm xếp hạng theo từ khóa chuẩn của truy xuất thông tin~\cite{robertson2009}. Điểm của đoạn $c$ đối với câu hỏi $q$ là tổng theo từng từ $t$ của $q$:
\begin{equation}
\mathrm{BM25}(q,c)=\sum_{t\in q}\mathrm{idf}(t)\cdot\frac{f(t,c)\,(k_1+1)}{f(t,c)+k_1\left(1-b+b\,\frac{|c|}{\mathrm{avgdl}}\right)},
\label{eq:bm25}
\end{equation}
\begin{equation}
\mathrm{idf}(t)=\ln\!\left(1+\frac{N-n(t)+0{,}5}{n(t)+0{,}5}\right).
\end{equation}
Ký hiệu $f(t,c)$ chỉ số lần từ $t$ xuất hiện trong đoạn $c$, $|c|$ là độ dài đoạn, $\mathrm{avgdl}$ là độ dài trung bình của các đoạn, $N$ là số đoạn và $n(t)$ là số đoạn chứa $t$. Tham số $k_1$ điều khiển độ bão hòa của tần suất và $b$ điều khiển chuẩn hóa độ dài; đồ án dùng $k_1=1{,}2$ và $b=0{,}75$. BM25 mạnh với thuật ngữ và con số chính xác, là hai loại thông tin chiếm phần lớn trong nội dung Toán.

Truy xuất dày so cosin giữa vector của câu hỏi và vector của đoạn~\cite{karpukhin2020}. Cách này bắt được những câu hỏi diễn đạt khác với tài liệu, chẳng hạn trẻ hỏi ``số có phẩy là gì'' trong khi tài liệu viết ``số thập phân''. Điểm yếu của truy xuất dày nằm ở chỗ nó kém nhạy với số và ký hiệu chính xác.

\subsection{Truy xuất lai, hợp nhất thứ hạng và xếp hạng lại}
Truy xuất lai kết hợp nhiều bảng xếp hạng thành một. Phương pháp hợp nhất theo hạng nghịch đảo (RRF) cho điểm của một đoạn bằng tổng nghịch đảo thứ hạng của đoạn đó trong từng bảng~\cite{cormack2009}:
\begin{equation}
\mathrm{RRF}(c)=\sum_{r\in R}\frac{1}{k+\mathrm{rank}_r(c)},
\label{eq:rrf}
\end{equation}
với $R$ là tập các bảng xếp hạng và $k$ thường bằng $60$. RRF không cần chuẩn hóa điểm giữa các hệ khác nhau. Bộ kiểm chuẩn BEIR cho thấy không hệ truy xuất nào thắng ở mọi miền và truy xuất từ khóa thường rất mạnh với văn bản chuyên ngành~\cite{thakur2021}, nên kết hợp hai phương pháp là lựa chọn an toàn. Hệ thống xếp hạng lại các ứng viên theo tiêu chí của bài toán sau bước hợp nhất. Hình~\ref{fig:hybrid-retrieval} mô tả chuỗi này, trong đó tiêu chí xếp hạng lại của đồ án là phạm vi trên đồ thị tri thức và mục đích sư phạm.

\begin{figure}[ht]
\centering
\begin{tikzpicture}[font=\footnotesize,
  box/.style={draw=black,line width=0.5pt,fill=white,align=center,minimum height=11mm,text width=24mm,inner sep=2pt},
  gbox/.style={box,fill=gray!14},
  arr/.style={-{Stealth[length=2mm,width=1.6mm]},line width=0.6pt},
  lab/.style={font=\scriptsize\itshape,fill=white,inner sep=1pt}]
\node[box] (q) at (0,0) {\textbf{Câu hỏi}\\\scriptsize của học sinh};
\node[box] (bm) at (3.6,1.0) {\textbf{BM25}\\\scriptsize khớp từ khóa};
\node[box] (vec) at (3.6,-1.0) {\textbf{Vector nhúng}\\\scriptsize khớp ngữ nghĩa};
\node[box] (rrf) at (7.2,0) {\textbf{Hợp nhất RRF}\\\scriptsize theo thứ hạng};
\node[gbox] (rr) at (10.8,0) {\textbf{Xếp hạng lại}\\\scriptsize phạm vi theo đồ thị};
\node[box,draw=none,fill=none] (out) at (14.1,0) {\textbf{Top-$k$}\\\scriptsize đoạn nguồn};
\draw[arr] (q.east) -- ++(0.4,0) |- (bm.west);
\draw[arr] (q.east) -- ++(0.4,0) |- (vec.west);
\draw[arr] (bm.east) -- ++(0.4,0) |- ([yshift=2mm]rrf.west);
\draw[arr] (vec.east) -- ++(0.4,0) |- ([yshift=-2mm]rrf.west);
\draw[arr] (rrf) -- (rr);
\draw[arr] (rr) -- (out);
\end{tikzpicture}

\caption{Chuỗi truy xuất lai: hai bảng xếp hạng độc lập được hợp nhất bằng RRF, sau đó các ứng viên được xếp hạng lại theo phạm vi trên đồ thị tri thức.}
\label{fig:hybrid-retrieval}
\end{figure}

\subsection{Độ bám nguồn và chỉ số đánh giá truy xuất}
Độ bám nguồn đòi hỏi câu trả lời chỉ khẳng định điều có trong đoạn nguồn. Các hướng đo gồm thước đo quy gán~\cite{rashkin2023}, bộ chỉ số RAGAs~\cite{es2024} và sinh văn bản kèm trích dẫn~\cite{gao2023alce}. GraphRAG dùng đồ thị để tổng hợp thông tin trên cả kho tài liệu~\cite{edge2024}, còn đồ án dùng đồ thị để giới hạn phạm vi truy xuất.

Chất lượng truy xuất đo bằng hai chỉ số. Chỉ số Hit@$k$ của một câu hỏi có đoạn đúng đầu tiên ở hạng $r$ bằng $1$ nếu $r\le k$ và bằng $0$ nếu ngược lại. Hạng nghịch đảo của câu hỏi đó là $1/r$, và MRR là trung bình của hạng nghịch đảo trên mọi câu hỏi:
\begin{equation}
\mathrm{MRR}=\frac{1}{|Q|}\sum_{q\in Q}\frac{1}{r_q}.
\label{eq:mrr}
\end{equation}

\section{Đồ thị tri thức và mô hình người học}

\subsection{Đồ thị tri thức trong giáo dục}
Đồ thị tri thức biểu diễn các đơn vị kiến thức là đỉnh và quan hệ giữa chúng là cạnh. Quan hệ quan trọng nhất trong giáo dục là quan hệ tiên quyết: học sinh phải học $u$ trước $v$. Lý thuyết không gian tri thức mô hình hóa tập các trạng thái tri thức hợp lệ theo quan hệ tiên quyết~\cite{doignon1999}, và là nền của các hệ học thích ứng. Mô hình 2T-KT kết hợp mô hình ngôn ngữ với hai đồ thị, một đồ thị chung của môn học và một đồ thị riêng của từng học sinh, để mô phỏng cách giáo viên suy luận khi dữ liệu còn ít~\cite{li2026}. Đồ án kế thừa từ 2T-KT năm ý tưởng: đồ thị có kiểm chứng, ngưỡng tin cậy cho nhãn do mô hình đề xuất, khởi động nguội từ điểm của bài tiên quyết, kiểm khẳng định với bằng chứng và cách trình bày thí nghiệm loại bỏ từng thành phần.

Một chương học trong đồ án là một đồ thị có hướng $G=(V,E)$. Mỗi đỉnh là một tiết học, và mỗi cạnh mang một trong hai kiểu: tiên quyết (học $u$ trước $v$) hoặc vận dụng ($v$ vận dụng kiến thức của $u$). Đồ thị này đồng thời quyết định thứ tự dạy và phạm vi truy xuất, nên một đồ thị sai cấu trúc kéo theo cả hành vi dạy sai lẫn truy xuất sai.

\subsection{Truy vết tri thức}
Truy vết tri thức ước lượng học sinh đã nắm từng kỹ năng đến đâu từ lịch sử làm bài. Mô hình Bayes mô tả mỗi kỹ năng là một biến ẩn hai trạng thái, đã biết hay chưa biết, với bốn tham số: xác suất biết từ đầu, xác suất học được sau mỗi lần luyện, xác suất đoán đúng khi chưa biết và xác suất sơ suất khi đã biết~\cite{corbett1995}. Phân tích yếu tố hiệu năng dùng hồi quy logistic trên số lần đúng và sai~\cite{pavlik2009}. Truy vết tri thức sâu dùng mạng nơ-ron hồi quy học trực tiếp từ chuỗi bài làm~\cite{piech2015}. Pelánek tổng hợp sự đánh đổi giữa độ chính xác dự đoán và khả năng diễn giải của các mô hình này~\cite{pelanek2017}. Bảng~\ref{tab:kt} so sánh ba hướng đó với cách tiếp cận của đồ án.

\begin{table}[ht]
\centering
\caption{So sánh các phương pháp truy vết tri thức với cách tiếp cận của đồ án.}
\label{tab:kt}
\small
\begin{tabularx}{\linewidth}{@{}p{2.6cm}XXX@{}}
\toprule
Phương pháp & Biểu diễn & Dữ liệu cần để dùng & Khả năng diễn giải \\
\midrule
BKT & Biến ẩn hai trạng thái, bốn tham số & Phải ước lượng tham số từ dữ liệu & Cao \\
PFA & Hồi quy logistic trên số lần đúng, sai & Phải huấn luyện từ dữ liệu & Khá \\
DKT & Mạng nơ-ron hồi quy & Rất nhiều chuỗi bài làm & Thấp \\
Beta--Bernoulli (đồ án) & Phân phối Beta cho mỗi tiết và mỗi chiều & Không cần huấn luyện; tham số do khung đánh giá quy định & Cao; mỗi điểm truy ngược về một sự kiện \\
\bottomrule
\end{tabularx}
\end{table}

Đồ án chọn khả năng diễn giải và khả năng phát lại thay vì độ chính xác dự đoán, vì hệ thống chưa có dữ liệu học sinh thật đủ lớn để huấn luyện và vì giáo viên cần nhìn được bằng chứng đằng sau mỗi con số.

\subsection{Cập nhật Beta--Bernoulli}\label{sec:beta-ref}
Phân phối $\mathrm{Beta}(\alpha,\beta)$ là tiên nghiệm liên hợp của xác suất thành công trong mô hình Bernoulli. Phép cập nhật cho một bằng chứng có điểm $s\in[0,1]$ và trọng số $w$ là
\begin{equation}
\alpha\leftarrow\alpha+w\,s,\qquad \beta\leftarrow\beta+w\,(1-s),\qquad \mu=\frac{\alpha}{\alpha+\beta}.
\label{eq:beta}
\end{equation}
Phép cập nhật này có tính giao hoán: thứ tự của các bằng chứng không làm đổi kết quả cuối cùng, nên mức hiểu bài tính lại được từ sổ sự kiện theo bất kỳ thứ tự nào. Phép cập nhật cũng giải thích được, vì mỗi bằng chứng đóng góp một lượng rõ ràng vào $\alpha$ và $\beta$. Hình~\ref{fig:beta} minh họa cách phân phối thu hẹp và dịch chuyển khi bằng chứng tích lũy. Đồ án giữ một phân phối Beta cho mỗi tiết và mỗi chiều của khung đánh giá.

\begin{figure}[ht]
\centering
\begin{tikzpicture}
\begin{axis}[width=0.86\linewidth,height=6.2cm,domain=0.001:0.999,samples=120,
  xlabel={xác suất đã nắm tiết học $p$},ylabel={mật độ},xmin=0,xmax=1,ymin=0,ymax=4.2,
  legend style={font=\footnotesize,at={(0.03,0.97)},anchor=north west,draw=none,fill=none},
  tick label style={font=\footnotesize},label style={font=\small},ymajorgrids,
  axis lines=left]
\addplot[black,line width=0.9pt,dash pattern=on 3pt off 2pt]{1};
\addplot[black!60,line width=0.9pt]{12*x^2*(1-x)};
\addplot[black,line width=1.1pt]{252*x^6*(1-x)^2};
\legend{{$\mathrm{Beta}(1,1)$ chưa có bằng chứng},{$\mathrm{Beta}(3,2)$ sau vài bằng chứng},{$\mathrm{Beta}(7,3)$ sau nhiều bằng chứng}}
\end{axis}
\end{tikzpicture}

\caption{Phân phối Beta của xác suất đã nắm tiết học qua ba trạng thái bằng chứng. Phân phối càng thu hẹp thì ước lượng càng chắc chắn.}
\label{fig:beta}
\end{figure}

\section{Bandit ngữ cảnh và đánh giá ngoài chính sách}

\subsection{Bài toán bandit và đánh đổi giữa khai thác và thăm dò}
Bài toán bandit nhiều tay mô tả một tác nhân phải chọn lặp đi lặp lại một trong số các lựa chọn, gọi là cánh tay, và chỉ quan sát được phần thưởng của cánh tay đã chọn~\cite{auer2002}. Tác nhân không biết trước cánh tay nào tốt nhất. Mỗi lần chọn cánh tay đang có vẻ tốt nhất là khai thác kiến thức hiện có, còn mỗi lần thử một cánh tay ít được biết đến là thăm dò để tìm ra lựa chọn tốt hơn. Đánh đổi giữa hai việc này là vấn đề trung tâm của bài toán, và thước đo chất lượng của một thuật toán là hối tiếc tích lũy, tức tổng chênh lệch giữa phần thưởng của lựa chọn tốt nhất và phần thưởng của lựa chọn thuật toán đã chọn~\cite{lattimore2020}.

Dạy học là một bối cảnh tự nhiên của bài toán này. Mỗi lượt dạy đặt gia sư trước vài nước đi sư phạm cùng hợp lệ, chẳng hạn hỏi lại cách làm hoặc đưa gợi ý, và gia sư chỉ biết kết quả của nước đã chơi. Clement và cộng sự dùng bandit nhiều tay để chọn hoạt động cho từng học sinh trong một hệ gia sư thông minh~\cite{clement2015}, và Lan và Baraniuk đề xuất khung bandit ngữ cảnh để chọn hành động học tập cá nhân hóa~\cite{lan2016}.

\subsection{Bandit ngữ cảnh}
Bandit ngữ cảnh bổ sung vào bài toán một vector ngữ cảnh $x_t$ quan sát được trước mỗi lượt. Tác nhân nhận ở lượt $t$ ngữ cảnh $x_t$ và một tập cánh tay cho phép $\mathcal{A}_t$, chọn $a_t\in\mathcal{A}_t$ và nhận phần thưởng $r_t$ chỉ của cánh tay đã chọn. Mục tiêu là học chính sách $\pi(x,\mathcal{A})$ làm tổng phần thưởng lớn nhất, và hối tiếc sau $T$ lượt là
\begin{equation}
R_T=\sum_{t=1}^{T}\Big(\max_{a\in\mathcal{A}_t}\mu(x_t,a)-\mu(x_t,a_t)\Big),
\label{eq:regret}
\end{equation}
trong đó $\mu(x,a)=\mathbb{E}[r\mid x,a]$ là phần thưởng kỳ vọng. Ngữ cảnh cho phép lựa chọn tốt nhất thay đổi theo từng học sinh: nước đi hợp với một em đang tự suy nghĩ có thể không hợp với một em đang bế tắc.

\subsection{Thuật toán LinUCB và thăm dò $\varepsilon$-tham lam}
LinUCB giả định phần thưởng kỳ vọng của mỗi cánh tay là hàm tuyến tính của ngữ cảnh, $\mu(x,a)=\theta_a^{\top}x$, và ước lượng $\theta_a$ bằng hồi quy ridge trên các lượt đã chọn cánh tay đó~\cite{li2010linucb}. Thuật toán tính $\hat\theta_a=A_a^{-1}b_a$, với $A_a=\lambda I+\sum x\,x^{\top}$ và $b_a=\sum r\,x$ cộng trên các lượt chọn $a$, rồi chọn cánh tay có cận trên tin cậy lớn nhất:
\begin{equation}
a_t=\arg\max_{a\in\mathcal{A}_t}\ \hat\theta_a^{\top}x_t+\alpha\sqrt{x_t^{\top}A_a^{-1}x_t}.
\label{eq:linucb}
\end{equation}
Số hạng thứ hai là phần thưởng cho sự bất định: một cánh tay ít được chọn ở ngữ cảnh gần $x_t$ có $A_a^{-1}$ lớn nên được ưu tiên thăm dò. Hệ số $\alpha$ điều khiển mức thăm dò. Các tổng $\sum xx^{\top}$, $\sum rx$ và số lượt đều cộng được, nên trạng thái của LinUCB không phụ thuộc thứ tự các lượt chọn và tính lại được từ nhật ký bằng một phép gập, giống cách tính mức hiểu bài ở Mục~\ref{sec:beta-ref}. Đồ án thêm thăm dò $\varepsilon$-tham lam: với xác suất $1-\varepsilon$ thuật toán chọn theo công thức~\eqref{eq:linucb}, với xác suất $\varepsilon$ chọn ngẫu nhiên đều trong $\mathcal{A}_t$. Cách này làm mọi cánh tay có xác suất được chọn lớn hơn không và cho phép ghi lại xác suất của từng lượt chọn.

\subsection{Đánh giá ngoài chính sách}
Đánh giá một chính sách chưa từng chạy trên người học thật cần phương pháp ngoài chính sách: chỉ dùng nhật ký của một chính sách đã chạy. Ước lượng cho điểm theo nghịch đảo xác suất (IPS) dựa trên phép lấy mẫu của Horvitz và Thompson~\cite{horvitz1952}. Cho nhật ký $\{(x_t,a_t,r_t,p_t)\}$ với $p_t$ là xác suất chính sách ghi nhật ký đã chọn $a_t$, giá trị của chính sách đích tất định $\pi$ được ước lượng bởi
\begin{equation}
\hat V(\pi)=\frac{1}{T}\sum_{t=1}^{T}\frac{\mathbb{1}[\pi(x_t,\mathcal{A}_t)=a_t]\,r_t}{p_t}.
\label{eq:ips}
\end{equation}
Ước lượng này không chệch khi $p_t>0$ cho mọi cánh tay mà chính sách đích có thể chọn, và phương sai tăng khi chính sách đích khác xa chính sách ghi nhật ký~\cite{li2011replay}. Đồ án vì vậy ghi xác suất của mỗi lượt chọn trong nhật ký sự kiện, điều mà một chính sách tham lam thuần túy không làm được.

\section{Xử lý tiếng nói và đọc tài liệu}

\subsection{Nhận dạng và tổng hợp tiếng nói}
Whisper là mô hình Transformer mã hóa--giải mã nhận dạng tiếng nói, nhận đầu vào là phổ log-Mel của âm thanh và được huấn luyện có giám sát yếu trên khoảng 680\,000 giờ âm thanh nhiều ngôn ngữ, trong đó có tiếng Việt~\cite{radford2023}. Đồ án dùng \texttt{whisper-large-v3} trên Groq cho nhận dạng, và dùng Web Speech API của trình duyệt làm dự phòng khi dịch vụ lỗi.

Tổng hợp tiếng nói nơ-ron biến văn bản thành âm thanh tự nhiên. Đồ án dùng VieNeu-TTS, một hệ tổng hợp tiếng Việt mã nguồn mở theo giấy phép Apache 2.0~\cite{vieneu}. Bản v3 Turbo chạy trên CPU qua ONNX Runtime, với giọng nữ miền Bắc vì gia sư xưng ``cô''. Đo trên Apple M1 Pro cho thấy từ câu thứ hai, 5,6 giây âm thanh được tạo trong 4,2 giây, tức nhanh hơn thời gian thực.

Chữ Toán phải được chuẩn hóa trước khi đọc, sao cho giọng nói đọc như giáo viên đọc: ``3,05'' thành ``ba phẩy không năm'', ``3/10'' thành ``ba phần mười'' và ``cm$^2$'' thành ``xăng-ti-mét vuông''. Luân phiên lượt nói cũng cần quy ước riêng: hệ thống coi học sinh nói xong sau 1,5 giây im lặng, và cho phép học sinh ngắt lời cô.

\subsection{Đọc tài liệu của giáo viên}
Tài liệu của giáo viên gồm các tệp PDF có lớp chữ và các trang quét hoặc chụp. Trang có lớp chữ được trích trực tiếp, còn trang ảnh được đọc bằng mô hình thị giác--ngôn ngữ, vốn đọc được cả công thức và hỗn số. Mô hình có thể đọc sai số, nên đồ án kiểm chứng bằng cách đòi mọi số đọc được phải khớp giữa các nguồn và đánh dấu chỗ lệch để giáo viên quyết định.

\section{Phương pháp đánh giá hệ dùng mô hình ngôn ngữ}

Giám khảo là mô hình ngôn ngữ chấm đầu ra mở theo tiêu chí, và tương quan của giám khảo với người chấm khá cao~\cite{zheng2023}. Giám khảo vẫn có thiên lệch, chẳng hạn thích câu dài, thích văn phong giống mình và bỏ sót lỗi cùng loại. Đồ án vì thế dùng giám khảo khác họ với mô hình gia sư và ưu tiên chỉ số tất định mỗi khi có thể.

Phần thống kê dùng khoảng Wilson cho tỉ lệ~\cite{wilson1927}, bootstrap cho khoảng tin cậy của chênh lệch theo cặp~\cite{efron1993} và kiểm định McNemar cho các cặp bất đồng~\cite{mcnemar1947}, theo khuyến nghị cho xử lý ngôn ngữ của Dror và cộng sự~\cite{dror2018}. Tham số chỉ được chỉnh trên tập chỉnh, còn con số báo cáo lấy từ tập giữ riêng viết sau và được đóng băng bằng thẻ git. Thí nghiệm loại bỏ từng thành phần đo đóng góp của thành phần đó vào kết quả chung.

\section{Kỹ thuật phần mềm nền}

Nhật ký sự kiện chỉ ghi thêm là kỹ thuật lưu trạng thái bằng chuỗi sự kiện thay vì ghi đè. Mức hiểu bài là một phép chiếu tất định của chuỗi đó: cùng chuỗi luôn cho cùng kết quả, và có thể tính lại khi đổi công thức. Vòng một câu hỏi là một máy trạng thái hữu hạn với các pha như chờ trả lời, chờ lý do, đang hỏi vì sao và đã khép. Chính sách dạy là một hệ luật có thứ tự: luật đầu tiên thỏa điều kiện quyết định nước đi và mã luật được ghi lại để giải thích.

Ba kỹ thuật còn lại phục vụ tính đúng đắn của dữ liệu. Điều khiển đồng thời lạc quan gắn mỗi buổi học với một số phiên bản trạng thái và chỉ cho ghi khi phiên bản khớp. Mô hình đa thuê bao gắn mọi dữ liệu người học với một định danh trường và bắt buộc mọi truy vấn lọc theo định danh đó, vì MongoDB không có bảo mật mức dòng. Lập trình phòng thủ với dịch vụ bên ngoài gồm chuỗi mô hình dự phòng, cơ chế thử lại và lưu đệm vector nhúng.

\section{Tổng kết chương}

Chương này đã trình bày các cơ sở lý thuyết của đồ án. Phương pháp Socratic, giàn giáo, tự giải thích và xung đột nhận thức cho các ràng buộc hành vi mà gia sư phải giữ. Năm giới hạn cố hữu của mô hình ngôn ngữ giải thích vì sao quyết định không được giao cho mô hình. Truy xuất lai, đồ thị tri thức, cập nhật Beta--Bernoulli và nhật ký sự kiện cung cấp các thành phần tất định thay thế cho những việc mà mô hình không được làm. Bandit ngữ cảnh cho phép chọn nước đi theo kết quả quan sát được mà vẫn nằm trong tập nước đi hợp lệ. Chương~\ref{chap:design} sử dụng các cơ sở này để phân tích yêu cầu và đề xuất kiến trúc hệ thống.

